% Avsnitt 17.7, ur lectures/L17/README.md.

\section{Sammanfattning}
\label{c17:sec:review}

\needspace{6\baselineskip}
\subsection*{Det här ska du kunna nu}
\begin{itemize}
  \item Implementera mottagarvägen i samma tillståndsmaskin som sänder.
  \item Förklara hur en nod avgör att den förlorat arbitreringen, och vad den gör då.
  \item Sätta ihop \code{rx_id}, \code{rx_dlc} och \code{rx_data} ur den mottagna bitströmmen och
    höja \code{rx_valid} vid rätt tillfälle.
  \item Köra systemtestbänken med två noder på en buss och läsa vad den rapporterar.
  \item Säga vad \code{error} betyder, vilka fyra orsaker som sätter den, och när den nollställs.
\end{itemize}

\needspace{6\baselineskip}
\subsection*{Frågor att testa dig själv med}
\begin{enumerate}
  \item Hur upptäcker en nod att den förlorat arbitreringen, och exakt vilka bitar jämför den?
  \item Vad gör en nod som förlorat arbitreringen, och varför är det inte ”försök igen omedelbart”?
  \item Varför måste \code{rx_shift_reg}s enable föregripas en cykel?
  \item En nod kan inte sända och ta emot samtidigt. Var i tillståndsmaskinen syns det?
  \item Vad sätter \code{error}, och vad nollställer den?
  \item Vad bevisar \code{can_controller_tb} som de fyra modultestbänkarna tillsammans inte gör?
\end{enumerate}

\needspace{6\baselineskip}
\subsection*{Riktvärde}
Här går projektet över från CAN till SPI. En grupp som inte fått \code{can_controller_tb} att passera
bör prioritera det framför registerbanken: SPI-halvan går att bygga och testa mot en kontroller som
ännu inte fungerar, men bring-upen i \chapref{c19:ch} gör det inte.

Systemtestbänken hoppas över av bygget tills kontrollern har en mottagarväg. Ser du fortfarande
”skipped” när du tycker att den borde köra, kör med \code{CI_BUILD_ALL=1}; bilaga~\ref{app:sim}
förklarar varför den kontrollen fungerar så.

\needspace{6\baselineskip}
\subsection*{Blicken framåt}
\Chapref{c18:ch} tar registerbanken, och SPI från vågformen och uppåt. Här börjar den halva som gör
kontrollern nåbar från en mikrokontroller.
